% TeX root = ../Main.tex

% First argument to \section is the title that will go in the table of contents. Second argument is the title that will be printed on the page.
\section[Question -- {\it Sobel filter vs Canny edge detector}]{}

\subsection{Sobel filter}
The Sobel filter is a efficient approximation of the first gaussian derivative (used to compute the gradients
in an image). It is a separable filter composed by a 1D kernel derivation x 1D kernel smoothing. 

\subsection{Canny edge detector}
The canny edge detector is a edge detector algorithm, which involves the Sobel filter for the gradient computation, but the
algorithm does not ends up finding the gradient but analyzes all the edges found by the use of the filter.
Indeed, the proper algorithm is composed of these steps:
\begin{itemize}
    \item Gaussian smoothing, since the derivatives are not robust to gaussian noise, to smooth the image
    \item Sobel filter for gradient computation
    \item NMS to filter all the non maxima gradient pixels in ad edge, ensuring the edges to be 1 pixel thick
    \item Linking and Hyisteresis thresholding, all the pixels with weak gradient are kept if connected to a pixel with a strong gradient
\end{itemize}
